\documentclass[12pt,oneside,openany]{book}

% Metadata - populated from workspace.yml
\newcommand{\PDFTitle}{DATA is DNA}
\newcommand{\PDFTitleFront}{DATA is DNA}
\newcommand{\PDFCoverTitle}{DATA is DNA}
\newcommand{\PDFSubject}{}
\newcommand{\PDFKeywords}{}
\newcommand{\PDFAuthor}{m-mdy-m}
\newcommand{\PDFAuthorEmail}{bitsgenix@gmail.com}
\newcommand{\PDFURL}{https://github.com/papyrxis/data-is-dna}

% Import workspace preset (auto-generated from workspace.yml)
% Auto-generated by Papyrxis workspace sync
% DO NOT EDIT - This file is regenerated on each build
% Edit workspace.yml to configure components
% Generated: Wed May 27 16:00:44 +0330 2026

\usepackage{lmodern}
\usepackage[T1]{fontenc}
\usepackage[utf8]{inputenc}

\usepackage{charter}
\usepackage[scaled=0.92]{helvet}
\usepackage[scaled=0.85]{beramono}
\usepackage[charter]{mathdesign}
\usepackage[english]{babel}
\usepackage{csquotes}
\usepackage{lettrine}
\usepackage{microtype}
\usepackage{mathtools}
\usepackage{amsmath}
\usepackage{amssymb}
\usepackage{amsthm}
\usepackage{bm}

% Delimiters
\DeclarePairedDelimiter\floor{\lfloor}{\rfloor}
\DeclarePairedDelimiter\ceil{\lceil}{\rceil}
\DeclarePairedDelimiter\abs{\lvert}{\rvert}
\DeclarePairedDelimiter\norm{\lVert}{\rVert}
\DeclareMathOperator*{\argmin}{arg\,min}
\DeclareMathOperator*{\argmax}{arg\,max}
\DeclareMathOperator{\poly}{poly}
\DeclareMathOperator{\polylog}{polylog}
\DeclareMathOperator{\lcm}{lcm}
\usepackage{graphicx}
\usepackage{caption}
\usepackage{subcaption}
\usepackage{tikz}
\usetikzlibrary{positioning,shapes.geometric,arrows,calc}
\usepackage{pdfpages}
\usepackage{array}
\usepackage{multicol}
\usepackage{booktabs}
\usepackage{hyperref}
\hypersetup{
  pdfauthor={Author},
  pdftitle={Title},
  pdfsubject={Subject},
  pdfkeywords={Keywords}
}
\usepackage{xcolor}
\usepackage{xstring}

\definecolor{arlizblue}{RGB}{0, 86, 179}         
\definecolor{arlizlightblue}{RGB}{51, 122, 204} 
\definecolor{arliznavy}{RGB}{25, 45, 85}  
\definecolor{arlizgray}{RGB}{95, 99, 104}  

\definecolor{headercolor}{RGB}{25, 45, 85} 
\definecolor{headerrule}{RGB}{0, 86, 179}    
\definecolor{footercolor}{RGB}{95, 99, 104} 
\definecolor{footerrule}{RGB}{180, 185, 190}

\definecolor{codebg}{RGB}{248, 249, 250} 

\definecolor{textprimary}{RGB}{33, 37, 41} 
\definecolor{textsecondary}{RGB}{108, 117, 125}

\colorlet{myblue}{arlizblue}
\colorlet{lightgray}{codebg}
\colorlet{headerblue}{headercolor}
\colorlet{footergray}{footercolor}

\definecolor{primary}{RGB}{30, 40, 50}
\definecolor{accent}{RGB}{80, 120, 160}
\definecolor{subtle}{RGB}{240, 242, 245}
\definecolor{border}{RGB}{200, 205, 210}
\definecolor{textgray}{RGB}{60, 60, 60}
\definecolor{codebackground}{RGB}{248, 248, 250} 
\usepackage[
  a4paper,
  top=2.5cm,
  bottom=2.5cm,
  inner=3cm,
  outer=2.5cm,
  headheight=1.2cm,
  headsep=0.8cm,
  footskip=1cm,
  marginparwidth=2cm,
  marginparsep=0.3cm
]{geometry}

\usepackage{setspace}
\setstretch{1.25}

\setlength{\parindent}{1.5em}
\setlength{\parskip}{0.4ex plus 0.2ex minus 0.1ex}

\tolerance=2000
\hbadness=1000
\emergencystretch=15pt

\usepackage{float}
\usepackage{placeins}
\usepackage{afterpage}

\renewcommand{\topfraction}{0.85}
\renewcommand{\bottomfraction}{0.70}
\renewcommand{\textfraction}{0.1}
\renewcommand{\floatpagefraction}{0.8}

\setcounter{topnumber}{3}
\setcounter{bottomnumber}{2}
\setcounter{totalnumber}{5}

\usepackage{enumitem}

\setlist[itemize]{
  leftmargin=20pt,
  itemsep=4pt plus 1pt minus 1pt,
  parsep=2pt plus 1pt minus 1pt,
  topsep=8pt plus 2pt minus 2pt
}

\setlist[enumerate]{
  leftmargin=25pt,
  itemsep=4pt plus 1pt minus 1pt,
  parsep=2pt plus 1pt minus 1pt,
  topsep=8pt plus 2pt minus 2pt
}

\setlist[description]{
  leftmargin=0pt,
  itemsep=6pt plus 2pt minus 1pt,
  parsep=2pt plus 1pt minus 1pt,
  topsep=8pt plus 2pt minus 2pt,
  font=\bfseries\color{accent}
}

\widowpenalty=10000
\clubpenalty=10000
\brokenpenalty=7500
\predisplaypenalty=10000
\postdisplaypenalty=2500
\displaywidowpenalty=3000

\lefthyphenmin=3
\righthyphenmin=3

\linepenalty=10
\adjdemerits=10000

\raggedbottom

\usepackage{tocloft}
\usepackage{titletoc}

\setlength{\cftbeforetoctitleskip}{0pt}
\setlength{\cftaftertoctitleskip}{3ex}
\setlength{\cftbeforepartskip}{2ex plus 0.5ex}
\setlength{\cftbeforechapskip}{1.2ex plus 0.3ex}
\setlength{\cftbeforesecskip}{0.8ex plus 0.2ex}
\setlength{\cftbeforesubsecskip}{0.4ex plus 0.1ex}

\setlength{\cftchapnumwidth}{3em}
\setlength{\cftsecnumwidth}{3em}
\setlength{\cftsubsecnumwidth}{3.5em}

\renewcommand{\cftpartfont}{\Large\sffamily\bfseries\color{primary}}
\renewcommand{\cftpartpagefont}{\Large\sffamily\bfseries\color{primary}}
\renewcommand{\cftpartleader}{\cftdotfill{\cftpartdotsep}}
\renewcommand{\cftpartdotsep}{4.5}

\renewcommand{\cftchapfont}{\large\sffamily\bfseries\color{primary}}
\renewcommand{\cftchappagefont}{\large\sffamily\bfseries\color{primary}}
\renewcommand{\cftchapleader}{\cftdotfill{\cftchapdotsep}}
\renewcommand{\cftchapdotsep}{2.6}

\renewcommand{\cftsecfont}{\sffamily\color{textprimary}}
\renewcommand{\cftsecpagefont}{\sffamily\color{textprimary}}
\renewcommand{\cftsecleader}{\cftdotfill{\cftsecdotsep}}

\renewcommand{\cftsubsecfont}{\small\sffamily\color{textsecondary}}
\renewcommand{\cftsubsecpagefont}{\small\sffamily\color{textsecondary}}
\usepackage{titlesec}

\titlespacing*{\chapter}
  {0pt}
  {0pt}
  {2.5ex plus 0.5ex}

\titlespacing*{\section}
  {0pt}
  {2ex plus 0.4ex minus 0.2ex}
  {1.5ex plus 0.3ex minus 0.1ex}

\titlespacing*{\subsection}
  {0pt}
  {1.5ex plus 0.3ex minus 0.1ex}
  {1ex plus 0.2ex minus 0.1ex}

\titlespacing*{\subsubsection}
  {0pt}
  {1ex plus 0.2ex minus 0.1ex}
  {0.8ex plus 0.1ex minus 0.1ex}
\usepackage{fancyhdr}
\usepackage{lastpage}
\usepackage{ifthen}

\fancyhf{}
\renewcommand{\headrulewidth}{0pt}
\renewcommand{\footrulewidth}{0pt}

\newcommand{\frontmatterpagenumber}{%
  \pagenumbering{roman}%
  \setcounter{page}{1}%
}

\newcommand{\mainmatterpagenumber}{%
  \pagenumbering{arabic}%
  \setcounter{page}{1}%
}

\newcommand{\headerrule}{%
  \vbox to 0pt{%
    \vss%
    \hbox to \headwidth{%
      \textcolor{accent}{\leaders\hrule height 1.2pt\hfill}%
    }%
    \vskip 0.2pt%
    \hbox to \headwidth{%
      \textcolor{accent!60}{\leaders\hrule height 0.4pt\hfill}%
    }%
    \vskip 2pt%
  }%
}

\newcommand{\footerrule}{%
  \vbox to 0pt{%
    \vskip -4pt%
    \hbox to \headwidth{%
      \textcolor{border}{\leaders\hrule height 0.6pt\hfill}%
    }%
    \vss%
  }%
}

\newcommand{\decorativerule}{%
  \textcolor{accent}{\rule{8pt}{0.8pt}}%
  \,\textcolor{accent!70}{\textbullet}\,%
  \textcolor{accent}{\rule{8pt}{0.8pt}}%
}

\newcommand{\headerfont}{\sffamily\small\bfseries}
\newcommand{\footerfont}{\sffamily\footnotesize}
\newcommand{\pagenumberfont}{\sffamily\small\bfseries}

\fancypagestyle{main}{%
  \fancyhf{}
  
  \fancyhead[RO]{%
    \parbox[b]{\textwidth}{%
      \raggedleft%
      \headerfont\color{primary}%
      {\nouppercase{\rightmark}}%
      \vskip 2pt%
    }%
  }
  
  \fancyhead[LO]{%
    \parbox[b]{\textwidth}{%
      \raggedright%
      \headerfont\color{accent}%
      \textsc{\PDFTitleFront}%
      \vskip 2pt%
    }%
  }
  
  \fancyhead[LE]{%
    \parbox[b]{\textwidth}{%
      \raggedright%
      \headerfont\color{primary}%
      {\nouppercase{\leftmark}}%
      \vskip 2pt%
    }%
  }
  
  \fancyhead[RE]{%
    \parbox[b]{\textwidth}{%
      \raggedleft%
      \headerfont\color{accent}%
      \textsc{\PDFTitleFront}%
      \vskip 2pt%
    }%
  }
  
  \fancyfoot[LE]{%
    \footerfont\color{accent}%
    \textbf{\large\thepage}%
    \textcolor{border}{\,\textbar\,}%
    \textcolor{textsecondary}{\pageref{LastPage}}%
  }
  
  \fancyfoot[RO]{%
    \footerfont\color{accent}%
    \textbf{\large\thepage}%
    \textcolor{border}{\,\textbar\,}%
    \textcolor{textsecondary}{\pageref{LastPage}}%
  }
  
  \fancyfoot[RE]{%
    \footerfont\color{textsecondary}%
    \textit{Open Source}%
  }
  
  \fancyfoot[LO]{%
    \footerfont\color{textsecondary}%
    \textit{\the\year}%
  }
  
  \renewcommand{\headrulewidth}{0.6pt}
  \renewcommand{\footrulewidth}{0.4pt}
  \renewcommand{\headrule}{\hbox to\headwidth{\color{accent}\leaders\hrule height \headrulewidth\hfill}}
  \renewcommand{\footrule}{\hbox to\headwidth{\color{border}\leaders\hrule height \footrulewidth\hfill}}
}

\fancypagestyle{chapter}{%
  \fancyhf{}
  \renewcommand{\headrulewidth}{0pt}
  \renewcommand{\footrulewidth}{0.4pt}
  
  \fancyfoot[C]{%
    \footerfont\color{accent}%
    \decorativerule\quad%
    \textbf{\large\thepage}%
    \quad\decorativerule%
  }
  
  \renewcommand{\footrule}{\hbox to\headwidth{\color{accent!60}\leaders\hrule height \footrulewidth\hfill}}
}

\fancypagestyle{frontmatter}{%
  \fancyhf{}
  
  \fancyhead[C]{%
    \headerfont\color{primary}%
    \textsc{\PDFTitleFront}%
  }
  
  \fancyfoot[C]{%
    \pagenumberfont\color{textsecondary}%
    \textbf{\thepage}%
  }
  
  \renewcommand{\headrulewidth}{0.5pt}
  \renewcommand{\footrulewidth}{0.3pt}
  \renewcommand{\headrule}{\color{primary!70}\hrule width\headwidth height\headrulewidth}
  \renewcommand{\footrule}{\color{textsecondary!50}\hrule width\headwidth height\footrulewidth}
}

\fancypagestyle{part}{%
  \fancyhf{}
  \renewcommand{\headrulewidth}{0pt}
  \renewcommand{\footrulewidth}{0pt}
  
  \fancyfoot[C]{%
    \footerfont\color{primary}%
    \decorativerule\quad%
    \textbf{\Large\thepage}%
    \quad\decorativerule%
  }
}

\fancypagestyle{backmatter}{%
  \fancyhf{}
  
  \fancyhead[LE,RO]{%
    \headerfont\color{primary}{\nouppercase{\leftmark}}%
  }
  
  \fancyhead[LO,RE]{%
    \headerfont\color{textsecondary}{\textsc{\PDFTitleFront}}%
  }
  
  \fancyfoot[LE,RO]{%
    \footerfont\color{primary}{\textbf{\thepage}}%
  }
  
  \fancyfoot[LO,RE]{%
    \footerfont\color{textsecondary}{\textit{Back Matter}}%
  }
  
  \renewcommand{\headrulewidth}{0.5pt}
  \renewcommand{\footrulewidth}{0.3pt}
  \renewcommand{\headrule}{\headerrule}
  \renewcommand{\footrule}{\footerrule}
}

\fancypagestyle{empty}{%
  \fancyhf{}
  \renewcommand{\headrulewidth}{0pt}
  \renewcommand{\footrulewidth}{0pt}
}

\pagestyle{main}

\let\oldfrontmatter\frontmatter
\renewcommand{\frontmatter}{%
  \oldfrontmatter%
  \pagestyle{frontmatter}%
}

\let\oldmainmatter\mainmatter
\renewcommand{\mainmatter}{%
  \oldmainmatter%
  \pagestyle{main}%
}

\let\oldbackmatter\backmatter
\renewcommand{\backmatter}{%
  \oldbackmatter%
  \pagestyle{backmatter}%
}

\newcommand{\mainstyle}{\pagestyle{main}}
\usepackage{tcolorbox}

\newenvironment{partintro}{%
    \vspace{3em}
    \begin{center}
    \begin{minipage}{0.80\textwidth}
        \setlength{\parindent}{0pt}
        \setlength{\parskip}{1em}
        \itshape
        \color{textgray}
        \fontsize{11pt}{16pt}\selectfont
}{%
    \end{minipage}
    \end{center}
    \vspace{3em}
    \noindent
}

\newenvironment{chapterintro}{%
    \vspace{1.5em}
    \begin{tcolorbox}[
        enhanced,
        colback=white,
        colframe=white,
        boxrule=0pt,
        borderline west={4pt}{0pt}{arlizblue!60},
        sharp corners,
        left=16pt,
        right=12pt,
        top=12pt,
        bottom=12pt,
        breakable,
        before skip=12pt,
        after skip=18pt,
        parbox=false
    ]
    \setlength{\parindent}{0pt}
    \setlength{\parskip}{0.8em}
    \color{textgray}
    \fontsize{10.5pt}{15pt}\selectfont
}{%
    \end{tcolorbox}
    \noindent
}

\newenvironment{sectionintro}{%
    \vspace{1em}
    \begin{tcolorbox}[
        enhanced,
        colback=subtle,
        colframe=subtle,
        boxrule=0pt,
        sharp corners,
        left=12pt,
        right=12pt,
        top=10pt,
        bottom=10pt,
        breakable,
        before skip=10pt,
        after skip=14pt,
        parbox=false
    ]
    \setlength{\parindent}{0pt}
    \setlength{\parskip}{0.6em}
    \small
    \itshape
    \color{textgray}
}{%
    \end{tcolorbox}
    \noindent
}

\newenvironment{subsectionintro}{%
    \vspace{0.8em}
    \begin{quote}
    \setlength{\parindent}{0pt}
    \setlength{\parskip}{0.5em}
    \small
    \itshape
    \color{textgray}
}{%
    \end{quote}
    \vspace{0.5em}
    \noindent
}

\newenvironment{historical}{%
    \vspace{1em}
    \begin{tcolorbox}[
        enhanced,
        colback=arliznavy!2,
        colframe=arliznavy!30,
        boxrule=0.5pt,
        arc=3pt,
        left=14pt,
        right=14pt,
        top=12pt,
        bottom=12pt,
        breakable,
        before skip=10pt,
        after skip=14pt,
        title={\small\scshape\color{arliznavy!80}Historical Context},
        fonttitle=\sffamily\bfseries,
        coltitle=arliznavy!80,
        attach boxed title to top left={yshift=-3mm, xshift=8mm},
        boxed title style={
            colback=white,
            colframe=arliznavy!30,
            boxrule=0.5pt
        }
    ]
    \setlength{\parindent}{0pt}
    \setlength{\parskip}{0.7em}
    \small
    \color{textgray}
}{%
    \end{tcolorbox}
    \noindent
}

\newenvironment{example}[1][]{%
    \vspace{1em}
    \begin{tcolorbox}[
        enhanced,
        colback=green!3,
        colframe=green!60!black,
        boxrule=0.5pt,
        arc=2pt,
        left=12pt,
        right=12pt,
        top=10pt,
        bottom=10pt,
        breakable,
        before skip=10pt,
        after skip=12pt,
        title={\small\textbf{Example\if\relax\detokenize{#1}\relax\else: #1\fi}},
        fonttitle=\sffamily,
        coltitle=green!60!black
    ]
    \setlength{\parindent}{0pt}
    \setlength{\parskip}{0.6em}
}{%
    \end{tcolorbox}
    \noindent
}

\newenvironment{definition}[1][]{%
    \vspace{1em}
    \begin{tcolorbox}[
        enhanced,
        colback=arlizblue!5,
        colframe=arlizblue!70,
        boxrule=1pt,
        arc=0pt,
        sharp corners,
        left=14pt,
        right=14pt,
        top=12pt,
        bottom=12pt,
        breakable,
        before skip=10pt,
        after skip=14pt,
        title={\small\textbf{Definition\if\relax\detokenize{#1}\relax\else: #1\fi}},
        fonttitle=\sffamily\bfseries,
        coltitle=arlizblue!90
    ]
    \setlength{\parindent}{0pt}
    \setlength{\parskip}{0.6em}
}{%
    \end{tcolorbox}
    \noindent
}

\newenvironment{theorem}[1][]{%
    \vspace{1em}
    \begin{tcolorbox}[
        enhanced,
        colback=arliznavy!5,
        colframe=arliznavy!70,
        boxrule=1pt,
        arc=0pt,
        sharp corners,
        left=14pt,
        right=14pt,
        top=12pt,
        bottom=12pt,
        breakable,
        before skip=10pt,
        after skip=14pt,
        title={\small\textbf{Theorem\if\relax\detokenize{#1}\relax\else: #1\fi}},
        fonttitle=\sffamily\bfseries,
        coltitle=arliznavy!90
    ]
    \setlength{\parindent}{0pt}
    \setlength{\parskip}{0.6em}
    \itshape
}{%
    \end{tcolorbox}
    \noindent
}

\newenvironment{implementation}{%
    \vspace{1em}
    \begin{tcolorbox}[
        enhanced,
        colback=orange!5,
        colframe=orange!70!black,
        boxrule=0.5pt,
        arc=2pt,
        left=12pt,
        right=12pt,
        top=10pt,
        bottom=10pt,
        breakable,
        before skip=10pt,
        after skip=12pt,
        title={\small\texttt{Implementation Detail}},
        fonttitle=\sffamily\bfseries,
        coltitle=orange!70!black
    ]
    \setlength{\parindent}{0pt}
    \setlength{\parskip}{0.6em}
    \small
}{%
    \end{tcolorbox}
    \noindent
}

\usepackage{tocbibind}
\usepackage{imakeidx}
\makeindex
\usepackage[backend=biber,style=authoryear]{biblatex}
\usepackage{listings}
\usepackage{algorithm}
\usepackage{algpseudocode}

\lstset{
	basicstyle       = \ttfamily\small,
	backgroundcolor  = \color{codebg},
	frame            = leftline,
	framerule        = 2pt,
	rulecolor        = \color{arlizblue},
	xleftmargin      = 15pt,
	framexleftmargin = 15pt,
	breaklines       = true,
	breakatwhitespace= true,
	showstringspaces = false,
	showtabs         = false,
	tabsize          = 4,
	keepspaces       = true,
	columns          = flexible,
	numbers          = left,
	numberstyle      = \tiny\color{arlizgray}\ttfamily,
	numbersep        = 10pt,
	stepnumber       = 1,
	captionpos       = b,
	aboveskip        = 12pt,
	belowskip        = 12pt,
	escapeinside     = {(*@}{@*)},
	commentstyle     = \color{green!60!black}\itshape,
	keywordstyle     = \color{arlizblue}\bfseries,
	stringstyle      = \color{orange!80!black},
	identifierstyle  = \color{textprimary},
	morekeywords     = {array, structure, algorithm, complexity}
}
\algrenewcommand\algorithmicrequire{\textbf{Input:}}
\algrenewcommand\algorithmicensure{\textbf{Output:}}
\lstdefinestyle{pseudostyle}{
	basicstyle       = \ttfamily\small,
	mathescape       = true,
	keywords         = {if, then, else, elif, fi, while, do, od, for, to, return, function, procedure, end, begin},
	keywordstyle     = \color{arlizblue}\bfseries,
	commentstyle     = \color{green!60!black}\itshape,
	morecomment      = [l]{//},
	literate         = 
	{<=}{$\leq$}1
	{>=}{$\geq$}1
	{!=}{$\neq$}1
	{->}{$\rightarrow$}1
	{<-}{$\leftarrow$}1
	{:=}{$\gets$}1
	{=}{$=$}1
}
\usepackage{tcolorbox}
\tcbuselibrary{most}

\newtcolorbox{notebox}[1][]{%
  colback=blue!5!white,
  colframe=blue!75!black,
  title=Note,
  fonttitle=\bfseries,
  sharp corners,
  boxrule=1pt,
  left=8pt,
  right=8pt,
  top=8pt,
  bottom=8pt,
  breakable,
  #1
}

\newtcolorbox{tipbox}[1][]{%
  colback=green!5!white,
  colframe=green!75!black,
  title=Tip,
  fonttitle=\bfseries,
  sharp corners,
  boxrule=1pt,
  left=8pt,
  right=8pt,
  top=8pt,
  bottom=8pt,
  breakable,
  #1
}

\newtcolorbox{warningbox}[1][]{%
  colback=red!5!white,
  colframe=red!75!black,
  title=Warning,
  fonttitle=\bfseries,
  sharp corners,
  boxrule=1pt,
  left=8pt,
  right=8pt,
  top=8pt,
  bottom=8pt,
  breakable,
  #1
}

\newtcolorbox{examplebox}[1][]{%
  colback=yellow!5!white,
  colframe=yellow!75!black,
  title=Example,
  fonttitle=\bfseries,
  sharp corners,
  boxrule=1pt,
  left=8pt,
  right=8pt,
  top=8pt,
  bottom=8pt,
  breakable,
  #1
}

\newtcolorbox{definitionbox}[1][]{%
  colback=purple!5!white,
  colframe=purple!75!black,
  title=Definition,
  fonttitle=\bfseries,
  sharp corners,
  boxrule=1pt,
  left=8pt,
  right=8pt,
  top=8pt,
  bottom=8pt,
  breakable,
  #1
}
\makeatletter

\newcommand{\provideCommand}[1]{%
  \@ifundefined{#1}{}{\let#1\relax}%
}

\makeatother

% Complexity notation
\providecommand{\bigO}{\ensuremath{\mathcal{O}}}
\providecommand{\bigOmega}{\ensuremath{\Omega}}
\providecommand{\bigTheta}{\ensuremath{\Theta}}
\providecommand{\littleo}{\ensuremath{o}}
\providecommand{\littleomega}{\ensuremath{\omega}}
% Number sets
\providecommand{\Naturals}{\ensuremath{\mathbb{N}}}
\providecommand{\Integers}{\ensuremath{\mathbb{Z}}}
\providecommand{\Rationals}{\ensuremath{\mathbb{Q}}}
\providecommand{\Reals}{\ensuremath{\mathbb{R}}}
\providecommand{\Complex}{\ensuremath{\mathbb{C}}}

% Probability
\providecommand{\Prob}{\ensuremath{\mathbb{P}}}
\providecommand{\Expect}{\ensuremath{\mathbb{E}}}
\providecommand{\Var}{\ensuremath{\operatorname{Var}}}

% Set notation
\newcommand{\set}[1]{\{#1\}}
\newcommand{\setbuilder}[2]{\{#1 \mid #2\}}

% Text formatting
\newcommand{\keyword}[1]{\textcolor{accent}{\textbf{#1}}}
\newcommand{\term}[1]{\textit{#1}}

% Metadata
\hypersetup{
  pdftitle={Data Is DNA},
  pdfauthor={m-mdy-m},
  pdfsubject={},
  pdfkeywords={},
  colorlinks=true,
  linkcolor=accent,
  citecolor=accent,
  urlcolor=accent
}

% Version information
\newcommand{\ProjectVersion}{f282091}
\newcommand{\BuildDate}{2026-05-27 12:30:44 UTC}


\begin{document}

\frontmatter
\frontmatterpagenumber

% \begin{titlepage}
\centering

\begin{tikzpicture}[remember picture, overlay]
\shade[top color=accent!5, bottom color=accent!15, middle color=white] 
    (current page.south west) rectangle (current page.north east);

\foreach \x in {0,2,...,20} {
    \foreach \y in {0,3,...,27} {
        \node[circle, fill=accent!10, inner sep=0.8pt] 
            at ($(current page.south west) + (\x cm, \y cm)$) {};
    }
}

\draw[line width=6pt, color=primary, rounded corners=4pt] 
    ($(current page.north west) + (0.5cm, -0.5cm)$) -- 
    ($(current page.north west) + (3cm, -0.5cm)$) -- 
    ($(current page.north west) + (3cm, -3cm)$);

\draw[line width=6pt, color=primary, rounded corners=4pt] 
    ($(current page.south east) + (-0.5cm, 0.5cm)$) -- 
    ($(current page.south east) + (-3cm, 0.5cm)$) -- 
    ($(current page.south east) + (-3cm, 3cm)$);
\end{tikzpicture}

\vspace*{5cm}

{\fontsize{48}{54}\selectfont\bfseries\color{primary}
Data Is DNA
}

\vfill

{\large\color{accent}
m-mdy-m
}

\vspace{2cm}

\end{titlepage}

% Professional title page for Papyrxis workspace
\begin{titlepage}
\centering

% Background with subtle pattern
\begin{tikzpicture}[remember picture, overlay]
    % Gradient background
    \shade[top color=accent!3, bottom color=accent!10, middle color=white] 
        (current page.south west) rectangle (current page.north east);
    
    % Decorative corner elements
    \draw[line width=4pt, color=accent, rounded corners=3pt, opacity=0.7] 
        ($(current page.north west) + (0.5cm, -0.5cm)$) -- 
        ($(current page.north west) + (3cm, -0.5cm)$) -- 
        ($(current page.north west) + (3cm, -3cm)$);
    
    \draw[line width=4pt, color=accent, rounded corners=3pt, opacity=0.7] 
        ($(current page.south east) + (-0.5cm, 0.5cm)$) -- 
        ($(current page.south east) + (-3cm, 0.5cm)$) -- 
        ($(current page.south east) + (-3cm, 3cm)$);
    
    % Subtle dot pattern
    \foreach \x in {1,3,...,19} {
        \foreach \y in {2,5,...,26} {
            \node[circle, fill=accent!12, inner sep=0.6pt, opacity=0.4] 
                at ($(current page.south west) + (\x cm, \y cm)$) {};
        }
    }
\end{tikzpicture}

\vspace*{4cm}

% Main title with shadow effect
\begin{tikzpicture}[baseline=(title.base)]
    % Shadow
    \node[opacity=0.08, scale=1.05, text=black] (shadow) at (0.1,-0.1) {
        \fontsize{42}{48}\selectfont\bfseries\PDFTitle
    };
    
    % Main title
    \node (title) at (0,0) {
        \fontsize{42}{48}\selectfont\bfseries\textcolor{primary}{\PDFTitle}
    };
\end{tikzpicture}

\vspace{1.5cm}

% Subtitle (if exists)
\ifdef{\PDFSubtitle}{%
    \ifx\PDFSubtitle\empty\else
        {\Large\color{accent}\PDFSubtitle}
        \vspace{1cm}
    \fi
}{}

% Decorative separator
\begin{tikzpicture}
    \draw[line width=1.5pt, color=accent, opacity=0.6] (-4cm,0) -- (-1.5cm,0);
    \node[circle, fill=accent, inner sep=2pt, opacity=0.6] at (0,0) {};
    \draw[line width=1.5pt, color=accent, opacity=0.6] (1.5cm,0) -- (4cm,0);
\end{tikzpicture}

\vspace{2cm}

% Author information
{\large\textcolor{textsecondary}{\textit{\PDFAuthor}}}

\ifdef{\PDFAuthorEmail}{%
    \ifx\PDFAuthorEmail\empty\else
        \vspace{0.5cm}
        {\small\textcolor{textsecondary}{\PDFAuthorEmail}}
    \fi
}{}

\vfill

% Edition and version info
\begin{tikzpicture}
    \node[text width=10cm, align=center] (info) {
        \large\textsc{\textcolor{primary}{First Edition}}
    };
    
    \node[below=0.8cm of info, text=textsecondary] {
        \textit{Version \ProjectVersion}
    };
    
    % Decorative lines around edition
    \draw[line width=2pt, color=accent, opacity=0.4, rounded corners=1pt] 
        ($(info.west) + (-2cm, 0.2cm)$) -- ($(info.west) + (-0.2cm, 0.2cm)$);
    \draw[line width=2pt, color=accent, opacity=0.4, rounded corners=1pt] 
        ($(info.east) + (0.2cm, 0.2cm)$) -- ($(info.east) + (2cm, 0.2cm)$);
\end{tikzpicture}

\vspace{2cm}

% Footer with year and publisher info
\begin{tikzpicture}
    \node (footer) {
        \textcolor{textsecondary}{
            \textcopyright\ \the\year\ \PDFAuthor
        }
    };
    
    \ifdef{\PDFURL}{%
        \ifx\PDFURL\empty\else
            \node[below=0.4cm of footer, text=textsecondary, font=\small] {
                \url{\PDFURL}
            };
        \fi
    }{}
\end{tikzpicture}

\end{titlepage}

% Ensure next content starts on odd page (for two-sided printing)
\cleardoublepage
\thispagestyle{empty}

\vspace*{\fill}

\begin{flushleft}
\textbf{Data Is DNA}

\vspace{0.5em}

Written by m-mdy-m

\vspace{1em}

First Edition

\vspace{0.5em}

Published: 2026

\vspace{2em}

\textbf{Copyright Notice}

\vspace{0.5em}

Copyright © 2026 m-mdy-m

This work is licensed under the Creative Commons Attribution-ShareAlike 4.0 
International License. To view a copy of this license, visit:
\url{http://creativecommons.org/licenses/by-sa/4.0/}

\vspace{2em}

\textbf{Version Information}

\vspace{0.5em}

Version: \ProjectVersion \\
Built: \BuildDate

\vspace{1em}

For the latest version, visit: \\
\url{https://github.com/papyrxis/data-is-dna}

\vspace{1em}

\textit{Typeset with \LaTeX{} using Papyrxis workspace.}

\end{flushleft}

\clearpage

\tableofcontents
\chapter*{Preface}
\addcontentsline{toc}{chapter}{Preface}

\section*{Why Would Anyone Learn This Stuff?}

\lettrine{D}{ata is everywhere}. Every application you use, every image you store, every message you send, every scientific discovery, every financial transaction, every machine-learning model—each exists because information can be represented, organized, transformed, and interpreted as data.

Yet despite living in an age dominated by data, most programmers understand only fragments of what data truly is.

They learn syntax before representation. They manipulate variables without understanding memory. They use abstractions without examining the structures beneath them. Data types become keywords to memorize rather than mathematical and computational models with deep theoretical foundations.

This book exists because that superficial understanding is no longer sufficient.

Modern computing demands far more than the ability to write code that merely works. It demands understanding how information behaves inside machines, how memory stores meaning, how programming languages model reality through type systems, and how algorithms fundamentally depend upon the structure of data itself.

Without that understanding, programming becomes mechanical imitation.

With it, computing becomes intelligible.

\section*{The Name and Its Meaning}

\textbf{Data Is DNA} reflects a simple but profound idea: data is the genetic material of computation.

Just as biological DNA encodes the instructions that shape living organisms, data encodes the structures, states, behaviors, and meanings that shape computational systems. Every program, operating system, database, network protocol, and artificial intelligence model ultimately depends upon the representation and manipulation of data.

Data is not merely “information stored in memory.” It is the foundation upon which all software systems are constructed.

To understand computation deeply is therefore to understand data deeply.

\section*{What This Book Represents}

\textit{Data Is DNA} is not a gentle introduction to programming, nor is it a catalog of programming-language syntax. Instead, it represents a rigorous exploration of data from multiple perspectives: mathematical, computational, architectural, and philosophical.

The purpose of this work is not simply to teach how data types are used, but to explain why they exist, how they emerged, how they interact with hardware and software systems, and what theoretical principles govern their behavior.

The structure of the book reflects a deliberate progression through increasingly sophisticated layers of understanding:

\begin{itemize}
	\item \textbf{The Nature of Data} — What data fundamentally is, how humans abstract information, and how computation transforms meaning into representation
	
	\item \textbf{Fundamental Concepts of Data Types} — The essential principles underlying primitive and composite types in programming systems
	
	\item \textbf{Data Types in Formal Computer Science} — Mathematical models, formal systems, computability, and theoretical representations of information
	
	\item \textbf{Data Models and Abstractions in Programming} — How programming languages organize and abstract data through structures, objects, schemas, and models
	
	\item \textbf{Data Types and Algorithms} — The inseparable relationship between data representation and algorithmic behavior
	
	\item \textbf{Memory and Data Types} — How data physically exists inside computer memory, including layout, alignment, allocation, and representation
	
	\item \textbf{Type Theories in Modern Programming Languages} — Static typing, dynamic typing, dependent types, generics, polymorphism, and modern type systems
\end{itemize}

This progression is intentional. True understanding emerges only when abstract theory, practical implementation, and hardware reality are studied together rather than in isolation.

\section*{Who This Book Is For}

This book is written for readers who want more than surface-level familiarity with programming concepts.

It is intended for programmers, computer science students, systems engineers, language designers, and intellectually curious learners who want to understand what actually happens beneath the abstractions provided by modern software systems.

The material assumes mathematical maturity, patience, and willingness to engage deeply with complex ideas. Some chapters are theoretical. Others are highly technical. Several require careful and repeated study.

This is deliberate.

The foundations of computing cannot be understood through shortcuts.

\section*{A Living Work}

Like all serious technical works, this book evolves continuously. As better explanations emerge, as new connections become visible, as errors are discovered, and as computing itself evolves, the text will improve alongside that understanding.

Knowledge is not static, and neither is this work.

Your engagement—through criticism, corrections, discussion, experimentation, and implementation—contributes to that evolution.

\section*{A Note on Difficulty}

Some readers may initially find portions of this book challenging. Concepts such as formal type theory, memory representation, computational abstraction, and mathematical modeling demand sustained intellectual effort.

This difficulty is not accidental.

The modern software world often encourages shallow familiarity over deep comprehension. This book deliberately moves in the opposite direction. Its goal is not speed, but clarity. Not simplification, but understanding.

If certain sections feel difficult, persist. Re-read them. Experiment with the ideas. Implement the concepts. Deep understanding develops gradually through sustained engagement.

\section*{Begin}

Welcome to \textbf{Data Is DNA}.

Let us examine the foundations of computation through the language of data.

\vfill

\begin{quote}
\textit{``Information is the resolution of uncertainty.''}

\hfill--- \textsc{Claude Shannon}
\end{quote}

\clearpage
\input{frontmatter/introduction}

\chapter{The Nature of Information and Reality}

\section*{What Is Information?}
\section*{Data, Meaning, and Representation}
\section*{Patterns, Symbols, and Abstraction}
\section*{Information in Nature}
\section*{Entropy and Uncertainty}
\section*{Information Theory and Claude Shannon}
\section*{Signals, Noise, and Communication}
\section*{Analog vs Digital Reality}
\section*{Time, State, and Change}
\section*{The Relationship Between Information and Computation}

\vspace{1em}

This chapter establishes the philosophical and scientific foundations of the entire book. Before understanding data in computers, we must first understand information itself.

\chapter{Mathematical Foundations of Computation}

\section*{Numbers and Number Systems}
\section*{Binary as a Universal Representation}
\section*{Boolean Algebra}
\section*{Sets, Relations, and Functions}
\section*{Logic and Formal Reasoning}
\section*{Matrices, Vectors, and Structures}
\section*{Discrete Mathematics}
\section*{Probability and Information}
\section*{Graphs and Networks}
\section*{Mathematics as a Language of Representation}

\vspace{1em}

This chapter introduces the mathematical structures that make modern computation possible. These concepts later become the foundations of hardware logic, memory systems, algorithms, and type systems.

% =========================================================

\chapter{From Physics to Electronics}

\section*{Electricity and Electromagnetism}
\section*{Voltage, Current, Resistance, and Power}
\section*{Semiconductors and Silicon}
\section*{Transistors and Switching}
\section*{Digital Signals and Logic Levels}
\section*{Logic Gates and Circuits}
\section*{Clock Signals and Synchronization}
\section*{Integrated Circuits and Microchips}
\section*{Physical Limits of Computation}
\section*{Quantum Phenomena and Future Computing}

\vspace{1em}

This chapter bridges the gap between abstract mathematics and physical machines. Computation is not magic—it emerges from physical systems governed by electrical and quantum behavior.

% =========================================================

\chapter{Computer Architecture and Machine Organization}

\section*{What Is a Computer?}
\section*{The Von Neumann Model}
\section*{Processors and Execution Units}
\section*{Instruction Cycles}
\section*{Registers and CPU Architecture}
\section*{Memory Hierarchies}
\section*{Caches and Performance}
\section*{Input and Output Systems}
\section*{Storage Devices and Persistence}
\section*{Parallelism and Modern Architectures}

\vspace{1em}

This chapter explains how computers are organized internally and how hardware components cooperate to process information.

% =========================================================

\chapter{Memory, Representation, and Binary Data}

\section*{Bits, Bytes, and Binary Encoding}
\section*{Representing Numbers in Memory}
\section*{Floating-Point Representation}
\section*{Character Encoding Systems}
\section*{Images, Audio, and Multimedia Data}
\section*{Memory Layout and Addressing}
\section*{Pointers and References}
\section*{Alignment and Padding}
\section*{Endianness}
\section*{Compression and Data Transformation}

\vspace{1em}

This chapter examines how information physically exists inside machines and how raw binary patterns become meaningful computational structures.

% =========================================================

\chapter{Data Structures and Computational Organization}

\section*{Why Data Structures Exist}
\section*{Arrays and Sequential Memory}
\section*{Linked Structures}
\section*{Stacks and Queues}
\section*{Trees and Hierarchical Data}
\section*{Hash Tables and Associative Structures}
\section*{Graphs and Networks}
\section*{Sparse and Dense Representations}
\section*{Spatial and Multidimensional Structures}
\section*{Trade-offs Between Structures}

\vspace{1em}

This chapter introduces the fundamental organizational models used to structure and manipulate data efficiently.

% =========================================================

\chapter{Algorithms and Computational Processes}

\section*{What Is an Algorithm?}
\section*{Complexity and Efficiency}
\section*{Time and Space Analysis}
\section*{Searching and Sorting}
\section*{Recursion and Iteration}
\section*{Optimization Strategies}
\section*{Concurrency and Parallel Algorithms}
\section*{Probabilistic Computation}
\section*{Machine Learning and Data Processing}
\section*{Limits of Computability}

\vspace{1em}

This chapter demonstrates how algorithms emerge from the interaction between logic, memory, representation, and data structures.

% =========================================================

\chapter{Programming Languages and Type Systems}

\section*{Why Programming Languages Exist}
\section*{Syntax, Semantics, and Abstraction}
\section*{Primitive and Composite Types}
\section*{Static and Dynamic Typing}
\section*{Functions and Functional Models}
\section*{Objects and Data Abstraction}
\section*{Generic Programming and Polymorphism}
\section*{Type Safety and Verification}
\section*{Compilers, Interpreters, and Runtime Systems}
\section*{Modern Type Theory}

\vspace{1em}

This chapter connects low-level computational reality with high-level software abstractions and language design.

% =========================================================

\chapter{Systems, Networks, and Distributed Information}

\section*{Operating Systems and Resource Management}
\section*{Processes, Threads, and Scheduling}
\section*{Filesystems and Persistence}
\section*{Databases and Structured Information}
\section*{Computer Networks and Communication}
\section*{Distributed Systems}
\section*{Cloud Computing and Virtualization}
\section*{Security and Cryptography}
\section*{Fault Tolerance and Reliability}
\section*{The Internet as a Global Information System}

\vspace{1em}

This chapter expands computation beyond individual machines into interconnected systems operating at planetary scale.

% =========================================================

\chapter{Intelligence, Complexity, and the Future of Computation}

\section*{Artificial Intelligence and Information}
\section*{Neural Networks and Representation}
\section*{Complex Systems and Emergence}
\section*{Simulation and Computational Science}
\section*{Quantum Computing}
\section*{Biological Computation}
\section*{Self-Organizing Systems}
\section*{The Philosophy of Intelligence}
\section*{The Limits of Computation}
\section*{The Future of Data and Civilization}

\vspace{1em}

The final chapter examines the frontier of computation itself—where information, intelligence, physics, biology, and philosophy begin to converge.

%\input{backmatter/backmatter}
%\input{backmatter/note}

\printbibliography[title={Bibliography}]
\addcontentsline{toc}{chapter}{Bibliography}

\printindex
\addcontentsline{toc}{chapter}{Index}
\end{document}